% !TeX root = ../papers/nc-p3.tex
\section{同時零二次式の境界におけるHilbert関数の超過}
\label{sec:zero-pairing-defect}
節\ref{sec:klein-zero-quadratics}の表の左上、すなわち
$a=1$, $d=z$, $c=b=1/\lambda$を取る。開いたチャートの条件を保つ。
\begin{proposition}
この境界の$x=z-1$の二変数部分では全ての窓が非標準ASである。
また$(z,\lambda)=(3,2)$に固定した双一次曲線の全ての窓も非標準ASである。
この曲線では$x=4$の第0始点は経路長5まで標準次元を持つが、長さ6で92となる。
\end{proposition}
\begin{proof}
$H=0$を解くと、両行列は
\[
 A_0=\begin{pmatrix}1&1/\lambda\\1/\lambda&z\end{pmatrix},\qquad
 A_1=\begin{pmatrix}1/\lambda&1\\ z-x&1/\lambda+\lambda y\end{pmatrix},\qquad
 y=\frac{(z-1)^2(\lambda^2 x-(\lambda^2z-1))}
 {\lambda^2((\lambda^2z-1)x-(z-1)^2)}
\]
となる。$x=z-1$では分母の括弧が$(z-1)z(\lambda^2-1)$になる。
$\lambda^2=1$のときはこの制限上で分子も消えるため、
この二変数部分とは上式の有理グラフの閉包を指す。
$(3,2)$では分子と分母に共通根はなく、指定チャートの曲線全体がこのグラフに入る。
それぞれの関数体上で各始点の二次関係から語商を計算する。
前次数の基底に末尾の四矢を掛け、二つ前の基底に末尾の六関係を掛けた行を消去する。
これは前の位置までの関係を先に割り、最後の位置に入る関係だけを追加した商である。
関係基底の六列は一つの許容標本で独立な列を選ぶ。既知の両切断階数6により
関数体上で全関係を張る。計算は有理関数の多項式最大公約数による正規化と
厳密な掃き出しである。$\mathbb Q(z,\lambda)$上で$x=z-1$とすると
長さ4の四始点の次元は$(35,36,36,35)$、
$\mathbb Q(x)$上で$(z,\lambda)=(3,2)$とすると長さ6では$(92,84,85,84)$である。
後者の$x=4$で第0始点の長さ0から6までの次元は$(1,4,10,20,35,56,92)$である。
\path{computations/verify_zero_pairing_defect.py}がこれらを検算する。
階数6を保つ部分で関係空間は部分束なので、基底の取り替えを許せば、
経路次元は任意の特殊化で減らない。従って標準値35または84の超過は
それぞれの全閉包のチャート内の各点へ引き継がれる。
\end{proof}
